\chapter{讨论与局限性}
\label{chap:discuss}

本章讨论实验结果的解读、系统局限性以及未来工作方向。

\section{关于100\%准确率的审慎解读}

测试集100\%准确率需审慎解读，不应简单视为模型已完美。可能的原因包括：

\begin{enumerate}
    \item \textbf{测试集规模有限}：90样本的统计波动较大，100\%与94\%的差异在统计上不显著。
    \item \textbf{类别区分度较高}：三类样本的健康数据数值差异明显（如血氧$<$90\% vs $\geq$95\%），模型容易学习到强判别特征。
    \item \textbf{半合成数据的标签相关性}：生理/用药数据按风险等级规则生成，可能存在与标签的强相关性，使模型``作弊''。
\end{enumerate}

为排除纯shortcut learning（捷径学习）的可能，本工作通过以下三点验证：

\begin{itemize}
    \item \textbf{独立测试集}：测试集90样本与训练集720样本严格分离，无数据泄露
    \item \textbf{端到端真实数据验证}：ESC-50真实音频6/6正确，证明模型在真实输入上有效
    \item \textbf{消融实验四模态均有贡献}：若模型仅依赖单一shortcut，屏蔽其他模态不应导致准确率下降
\end{itemize}

尽管如此，在更大规模真实数据集上准确率可能回落。100\%应理解为``在本数据集分布上表现良好''，而非``模型已完美''。

\section{局限性}

\subsection{数据规模有限}

720训练样本对8.5M参数的Full模型偏小，存在过拟合风险。尽管早停和正则化（Dropout、label smoothing、weight decay）缓解了过拟合，但在更大数据集上的泛化能力有待验证。

\subsection{视频特征基于骨架而非端到端}

MediaPipe骨架特征丢失了外观/场景信息（如室内环境、物品位置）。3D CNN（如VideoMAE）可能捕获更丰富的视觉模式，但受限于CPU算力未采用。骨架特征的优势在于可解释性（关键点轨迹直观）和轻量性（无需GPU推理）。

\subsection{生理/用药为半合成数据}

生理和用药数据按医学规则生成，与真实采集数据存在分布差异：

\begin{itemize}
    \item 真实生理数据有噪声、缺失、个体差异
    \item 真实用药行为更复杂（忘记服药后补服、自行调整剂量等）
\end{itemize}

半合成数据的规则化可能放大了生理/用药模态的区分度，导致消融实验中health重要性偏高（-32.22\%）。真实临床数据的验证是未来必要的工作。

\subsection{CPU训练限制}

CPU训练耗时较长，无法进行大规模超参数搜索和数据增强实验。GPU环境下可尝试更多架构变体（如不同注意力头数、层数）和训练策略（如Mixup、Focal Loss）。

\subsection{模态权重接近均匀}

App展示的全局模态权重接近25\%均分，未充分分化。这是因为模型在小数据集上倾向于``平均使用''所有模态。消融实验比参数权重更可靠地反映模态重要性——即使权重接近均匀，屏蔽某模态仍会导致准确率下降，说明该模态确实在决策中发挥作用。

\section{未来工作}

\begin{enumerate}
    \item \textbf{扩充数据集}：收集更多视频类别（如不同跌倒姿态、不同场景）和真实生理采集数据，提升模型泛化能力。
    
    \item \textbf{预训练编码器}：使用VideoMAE、AST等预训练模型替代手工特征，捕获更丰富的模态语义。
    
    \item \textbf{GPU训练}：在GPU上进行更大batch size训练和超参数搜索，探索更优模型配置。
    
    \item \textbf{视频端到端学习}：用3D CNN替代骨架特征，实现视频模态的端到端可微学习。
    
    \item \textbf{真实临床验证}：与医疗机构合作，在真实老年人监护数据上验证系统有效性。
    
    \item \textbf{时序建模}：当前系统对单次输入做风险评估，未来可引入时序建模（如LSTM/Transformer对历史风险评估序列建模），实现趋势分析和预警。
\end{enumerate}
